\chapter{Finite-State Transducers and Approximate Matching}
\label{ch:finite-state-transducers}

A recognizer answers whether an input belongs to a language.
A spelling corrector must do more: associate an observed input with one or more dictionary outputs and rank those alternatives.
Replacing one letter, inserting a missing letter, and ignoring an extra letter all consume and produce different amounts of text.
A finite-state transducer exposes those two streams explicitly.

\section{Label transitions with two sequences}
\label{sec:ch79-relation}
A \term{finite-state transducer}, or FST, has a finite state set, initial and final states, and edges labeled with an input symbol and an output symbol, either of which may be $\epsilon$.
A path concatenates its input labels into $x$ and output labels into $y$ and therefore realizes the pair $(x,y)$.
The transducer defines a relation $R\subseteq\Sigma^*\times\Delta^*$, not necessarily a function.
Several paths can realize the same pair, and one input can have several different outputs.

An identity dictionary transducer consumes a stored word and emits that same word, rejecting other strings.
A morphology transducer can relate a lexical analysis to a surface spelling; inversion swaps its input and output labels and reverses the direction of that relation, not the direction of graph edges.
A byte encoder can map one scalar to several output bytes using an expanded state path.
Table~\ref{tab:ch79-labels} distinguishes common arc interpretations.

\begin{table}[htbp]
\caption{Input and output epsilon labels have different meanings.}
\label{tab:ch79-labels}
\centering
\begin{tabular}{lll}
\toprule
Arc label & Consumes & Produces \\
\midrule
$a:a$ & input a & output a \\
$a:b$ & input a & output b \\
$a:\epsilon$ & input a & nothing \\
$\epsilon:b$ & nothing & output b \\
$\epsilon:\epsilon$ & nothing & nothing \\
\bottomrule
\end{tabular}
\end{table}

An \term{epsilon transition} is not an end marker.
It moves through the machine without consuming a symbol on the indicated tape.
Cycles containing only epsilon on the input can produce unbounded output before another input symbol is consumed.
Algorithms that enumerate all paths therefore need a finiteness condition or budget even for a finite graph.

\section{Compose relations, not arbitrary determinized machines}
\label{sec:ch79-compose}
If $T$ maps $x$ to intermediate strings $y$ and $U$ maps those $y$ to $z$, their composition maps $x$ to $z$ whenever some compatible $y$ exists.
Without intermediate epsilon labels, composition follows pairs of states and synchronizes matching output labels from $T$ with input labels from $U$.
Figure~\ref{fig:ch79-compose} separates this relation from a destructive rewrite pipeline.
Composition preserves all admitted alternatives until a later projection or ranking step selects among them.

\begin{figure}[htbp]
\centering
\begin{tikzpicture}[node distance=15mm]
\node[basic block] (x) {observed input};
\node[basic block,below=of x] (t) {edit relation};
\node[basic block,below=of t] (u) {dictionary identity relation};
\node[basic block,below=of u] (z) {accepted candidate outputs};
\draw[cfg edge] (x)--(t);\draw[cfg edge] (t)--(u);\draw[cfg edge] (u)--(z);
\end{tikzpicture}
\caption{Composition restricts edit outputs to strings accepted by the dictionary.}
\label{fig:ch79-compose}
\end{figure}

General weighted composition needs epsilon-handling filters so different interleavings do not spuriously duplicate paths and alter their combined weights~\cite{mohri-2009-weighted}.
Algorithm~\ref{alg:ch79-compose} deliberately states the epsilon-free intermediate-label case.
It costs $O(V+E)$ graph work for the reachable product it actually constructs, with matching-label indexing costs included in edge generation; the reachable product can approach the full Cartesian state product.
The spelling composition in the next section avoids a generic filter by specializing one component to a fixed input and the other to an epsilon-free trie.

\begin{algorithm}[htbp]
\caption{Construct reachable composition when intermediate labels contain no epsilon.}
\label{alg:ch79-compose}
\begin{pseudocode}
\Require transducers $T,U$ with nonempty intermediate labels and finite graphs
\Ensure reachable product implementing relational composition
\State enqueue initial state pairs; mark them discovered
\While{a discovered pair $(p,q)$ remains unexpanded}
\ForAll{$T$ arcs $(p,a,b,p')$ and $U$ arcs $(q,b,c,q')$}
\State emit arc from $(p,q)$ to $(p',q')$ with input $a$ and output $c$
\State discover and enqueue $(p',q')$ if new
\EndFor
\State mark pair final exactly when both component states are final
\EndWhile
\end{pseudocode}
\end{algorithm}

Projection onto one tape yields an acceptor for that tape's admitted strings.
It does not make the relation deterministic.
Unlike ordinary finite acceptors, arbitrary transducers cannot always be converted into equivalent deterministic sequential transducers.
Weighted determinization also requires suitable algebraic and structural conditions.
A compiler should not invoke an acceptor subset construction on paired labels and assume it has solved general transducer determinization.

\section{Represent edit distance as path weight}
\label{sec:ch79-edit}
The unit-cost edit transducer includes copy $a:a$ at cost zero, substitution $a:b$ for unequal symbols at cost one, deletion $a:\epsilon$ at cost one, and insertion $\epsilon:b$ at cost one.
We measure edits from observed query to candidate spelling.
Thus a missing query letter is inserted on the output tape.
Reversing that convention does not change ordinary symmetric Levenshtein distance but does matter for asymmetric learned costs.

In the \term{tropical semiring}, alternative paths combine by minimum, consecutive arc costs combine by addition, and unreachable weight is infinity.
The minimum weight realizing $(x,y)$ is its edit distance under these operations.
Probability models use different combination operations and must account for multiple paths rather than automatically keeping just the cheapest one~\cite{mohri-2009-weighted}.
Negative weights also invalidate a direct Dijkstra argument.

For query \code{ct}, producing \code{cat} copies c, inserts a, and copies t, for cost one.
Producing \code{coat} requires two inserted letters and costs two.
Table~\ref{tab:ch79-distance} gives the independent dynamic-programming rows for \code{ct} against \code{cat}.
Entry $D[i,j]$ is the cost to transform the first $i$ query symbols into the first $j$ candidate symbols.

\begin{table}[htbp]
\caption{Dynamic programming and transducer paths compute the same unit-cost distance.}
\label{tab:ch79-distance}
\centering
\begin{tabular}{lrrrr}
\toprule
Query prefix & empty & c & ca & cat \\
\midrule
empty & 0 & 1 & 2 & 3 \\
c & 1 & 0 & 1 & 2 \\
ct & 2 & 1 & 1 & 1 \\
\bottomrule
\end{tabular}
\end{table}

The recurrence takes the minimum of deletion $D[i-1,j]+1$, insertion $D[i,j-1]+1$, and match/substitution $D[i-1,j-1]+[x_i\ne y_j]$.
Filling the full table costs $O(mn)$ time; keeping only the previous row costs $O(n)$ storage if only the final distance is needed.
Recovering an edit script requires additional predecessor information or a reconstruction algorithm.
Adjacent transposition is not one of these primitive edits: \code{ab} to \code{ba} costs two, not one.

\section{Search the implicit dictionary product}
\label{sec:ch79-search}
Fix a query $x$ and dictionary trie state $q$.
The product state $(i,q)$ records consumed query length and emitted dictionary prefix.
Deletion advances only $i$; insertion follows only a trie edge; copy/substitution advances both.
Acceptance requires $i=|x|$ and a final trie state.
Figure~\ref{fig:ch79-product} shows the local transition shapes.

\begin{figure}[htbp]
\centering
\begin{tikzpicture}[node distance=16mm and 18mm]
\node[basic block] (s) {$(i,q)$};
\node[basic block,right=of s] (d) {$(i+1,q)$};
\node[basic block,below=of s] (in) {$(i,q')$};
\node[basic block,right=of in] (b) {$(i+1,q')$};
\draw[cfg edge] (s)--node[above,font=\small]{delete: 1}(d);
\draw[cfg edge] (s)--node[left,font=\small]{insert: 1}(in);
\draw[cfg edge] (s)--node[above,sloped,font=\small]{copy or substitute}(b);
\end{tikzpicture}
\caption{Every specialized edit transition advances the query, the trie, or both.}
\label{fig:ch79-product}
\end{figure}

Algorithm~\ref{alg:ch79-product} performs shortest-path search while discarding paths already beyond radius $k$.
Every dictionary state has a unique prefix, so keeping the best cost for $(i,q)$ cannot discard a different word reaching that same state.
This proof fails if suffix states are merged without retaining additional output-path information.

\begin{algorithm}[htbp]
\caption{Enumerate dictionary candidates through an implicit unit-edit FST product.}
\label{alg:ch79-product}
\begin{pseudocode}
\Require query $x$, trie, nonnegative radius $k$, expansion budget
\Ensure complete minimum-cost candidate set, or explicit budget exhaustion
\State distance of $(0,root)$ is zero; insert it into a min-priority queue
\While{queue is nonempty}
\State remove a minimum-cost state; ignore stale entries
\State charge expansion budget; abort without claiming completeness if exhausted
\State record final dictionary word if all query symbols are consumed
\State relax deletion, insertion, and copy/substitution transitions within $k$
\EndWhile
\State return distinct candidate words with their minimum costs
\end{pseudocode}
\end{algorithm}

For $V$ reached product states and $E$ generated transitions, heap search costs $O((V+E)\log(V+E))$ with a lazy-entry binary heap and $O(V+E)$ worst-case stored entries.
The full product has at most $(|x|+1)Q$ states for $Q$ trie states, though radius pruning often avoids much of it.
No universal small-search guarantee follows from a small radius when the dictionary has many nearby branches.
The next chapter adds deterministic candidate ranking and tested exhaustion behavior.

\section{Exercises}
\label{sec:ch79-exercises}
\begin{enumerate}
\item\label{ex:ch79-edits} Give edit paths from \code{ct} to \code{cat} and \code{coat}, naming the input and output labels.
\item\label{ex:ch79-transpose} Compute the distance from \code{ab} to \code{ba} and explain why it differs from a distance with unit transposition.
\item\label{ex:ch79-product} A trie has 31 states and the query has four symbols.
Bound the unpruned specialized product's state count.
\item\label{ex:ch79-semiring} Two paths realize the same pair with costs two and three.
Give their tropical combination and explain why summing them is a different model.
\end{enumerate}

\section{Further reading}
\label{sec:ch79-reading}
Mohri develops weighted composition, epsilon processing, determinization conditions, and shortest-distance algorithms~\cite{mohri-2009-weighted}.
The string-algorithm text supplies complementary exact and approximate matching foundations~\cite{crochemore-2007-strings}.
